\documentclass[10pt,a4paper]{amsart}
\usepackage[margin=19mm]{geometry}
\usepackage{amsmath,amssymb,mathtools}
\usepackage[scr=rsfs,cal=euler]{mathalfa}
\usepackage{xfrac}
\usepackage[unicode,hidelinks,bookmarks=false]{hyperref}
\newcommand{\ie}{i.e.\ }
\newcommand{\cf}{cf.\ }
\newcommand{\resp}{respectively\ }
\newcommand{\Subsection}{\subsection}
\providecommand{\nobreakdash}{\nobreak\hbox{-}}
\setlength{\emergencystretch}{2em}
\multlinegap0pt
\usepackage{amssymb}
\usepackage{stackengine}
\usepackage{xcolor}
\usepackage[all]{xy}
\usepackage{tikz}
\usepackage{enumerate}
\newtheorem{lemma}{Lemma}
\newcommand{\mc}[1]{\mathcal{#1}}
\title{Brauer group of moduli stacks of parabolic principal bundles over a curve}
\author{Indranil Biswas, Sujoy Chakraborty}
\date{Source: arXiv:2602.14579v1 (CC BY 4.0)}
\begin{document}
\maketitle

\noindent\emph{Source and licence.} Brauer group of moduli stacks of parabolic principal bundles over a curve. Version arXiv:2602.14579v1 (CC BY 4.0), distributed by arXiv under the Creative Commons Attribution 4.0 International licence, http://creativecommons.org/licenses/by/4.0/, as recorded in the arXiv licence metadata for this exact version and shown on the abstract page https://arxiv.org/abs/2602.14579v1. Full licence notice: \texttt{licenses/arxiv-2602.14579-CC-BY-4.0.txt}.\par\smallskip

\noindent\textit{Editorial scope.} Counted unit: the article's lemma lem:eigenbasis (source0.tex lines 309--316). The article's complete original proof of it is delivered with it (source0.tex lines 318--327, one \texttt{proof} environment of 10 lines). No mathematics is written, repaired or re-derived: the statement and the proof are the article's own lines, in the article's own notation, and no retained line is substituted or re-flowed.  No further source-local context is needed: every cross-reference of every retained line resolves inside the two delivered spans.  Nothing was omitted from any retained span, so the package carries no omission record.  No retained line contains a citation, so the wrapper carries no bibliography and no bibliography entry.  No retained line contains a figure environment, an image-inclusion command or drawing code, so no image is delivered and the package declares no asset dependency.  Editorial formatting only, in addition: an AMS \texttt{amsart} preamble that restates the article's own declarations and macros used by the retained lines, namely the environments \texttt{lemma} on separate counters, together with the standard packages those lines need.  The retained environments print in this document as Lemma 1 in order of appearance, whereas the article prints the same items with the numbering of its own full document.  The complete native source of the article as distributed in the arXiv e-print for this exact version is bundled unchanged as \texttt{sources/194721/source0.tex}.
	\begin{lemma}\label{lem:eigenbasis}
		Let $V$ be a finite-dimensional vector space over a field $k$ with a diagonalizable endomorphism
		$\varphi\,:\, V\,\longrightarrow \,V$ and a filtration by subspaces
		$$V\ =\ V_1\ \supsetneq\ V_2\ \supsetneq \ \cdots\ \supsetneq V_\ell \ \supsetneq\ 0$$
		such that $\varphi(V_i) = V_i$ for all $1\, \leq\, i\, \leq\, \ell$. There exists a basis $\mc{B}_i$ of $V_i$
		consisting of eigenvectors for $\varphi$, satisfying $\mc{B}_i\,\supset\, \mc{B}_{i+1}$ for all
		$1\,\leq\, i\,\leq \,\ell$.
	\end{lemma}

	\begin{proof}
		By standard theory of semisimple modules,  $\varphi|_{V_i}$ is diagonalizable on $V_i$ for each $i$, so
		there exists a basis $\mc{B}_{i}$ of $V_{i}$ consisting of eigenvectors. Moreover, one can write
		$V_{i-1}\,=\,W_{i-1}\oplus V_{i}$ for some subspace $W_{i-1}$ which is invariant under $\varphi$. Thus,
		one can start with a basis $\mc{B}_{\ell}$ of $V_{\ell}$ consisting of eigenvectors, next choose a direct
		summand $W_{\ell-1}$ of $V_{\ell}$ in $V_{\ell-1}$, namely $V_{\ell-1} \,=\, W_{\ell-1}\oplus V_{\ell}$, and choose an
		eigenbasis of $W_{\ell-1}$;  its union with $\mc{B}_{\ell}$ produces a basis $\mc{B}_{\ell-1}$ for $V_{\ell-1}$
		consisting of eigenvectors and satisfying $\mc{B}_{\ell-1}\,\supset\,\mc{B}_{\ell}$. Proceeding inductively, one
		can produce bases $\mc{B}_i$ of $V_i$. This proves the lemma.
	\end{proof}

\end{document}
